\section{Cronograma de Execução}
O experimento descrito neste capítulo é executado ao longo do semestre
2026.2, com início em 10 de agosto de 2026. A Tabela~\ref{tab:cronograma}
apresenta a distribuição das atividades ao longo do semestre.
 
\begin{table}[h]
\centering
\caption{Cronograma de execução do experimento — semestre 2026.2}
\label{tab:cronograma}
\begin{tabular}{p{2.6cm}p{9.5cm}}
\toprule
\textbf{Período} & \textbf{Atividade} \\
\midrule
10/08 -- 30/08 & Implementação do gerador de stream sintética e do módulo de
injeção de drift (quatro padrões temporais, severidade controlada via
Equação~\ref{eq:severidade}). \\
31/08 -- 20/09 & Implementação do HDDDM em regime de janela deslizante e
validação de que a severidade observada converge para a severidade real em
condições sem ruído. \\
21/09 -- 11/10 & Implementação e adaptação a stream dos três classificadores
(Naive Bayes, Random Forest adaptativo, XGBoost via atualização em
micro-lotes). \\
12/10 -- 25/10 & Implementação das quatro estratégias de retreinamento (reset
completo, janela deslizante, incremental, reativação de modelos
históricos). \\
26/10 -- 08/11 & Execução do experimento fatorial completo (180 condições,
múltiplas repetições por condição) e coleta das métricas de avaliação. \\
09/11 -- 15/11 & Execução do protocolo de validação da questão derivada:
comparação entre sinais observáveis e características reais registradas em
cada condição. \\
16/11 -- 29/11 & Análise estatística dos resultados e redação do capítulo de
Resultados. \\
30/11 -- 06/12 & Redação da Discussão e Conclusão; revisão geral do
texto. \\
07/12 -- 13/12 & Preparação de slides e ensaio de defesa. \\
A definir pela coordenação & Defesa do TCC2. \\
\bottomrule
\end{tabular}
\end{table}